We present a falsification-first, galaxy-domain framework for explaining rotation-curve anomalies without assuming particulate cold dark matter. Working within standard general relativity (GR), we define a response-sector stress--energy $T^{\mathrm{resp}}_{\mu\nu}$ as the minimal effective source required so that the weak-field metric potentials inferred from kinematics satisfy the Einstein equation with baryons. This definition is ontologically agnostic and is best read as an EFT/constitutive-closure candidate. Observed rotation curves define a model-independent residual, which we compress with a one-parameter boundary-layer closure activated near a baryonic transition radius. Applied to 175 galaxies from the SPARC database, this model yields very strong improvement over a baryons-only model for 92\% of systems and passes a Bayesian Information Criterion (BIC) overfitting test in 97\% of cases. Crucially, we demonstrate that this 1-parameter response model is strongly preferred by BIC over standard 2-parameter NFW and Burkert dark matter halo models in over 90\% of the SPARC sample. Further, 5-fold cross-validation shows the response model has significantly lower out-of-sample prediction error (mean RMSE = 3.1 km/s) than the NFW model (mean RMSE = 4.9 km/s). The baryonic Tully--Fisher relation is recovered as an empirical consequence (Spearman $r_s = 0.95$). These results establish the response sector as a parsimonious and predictively powerful alternative to standard halo models for describing galactic dynamics.

\paragraph{Keywords.}
Galaxy rotation curves; dark matter alternatives; dark matter candidates; dark matter identity; general relativity; gravitation; effective field theory; response-sector stress--energy; boundary-layer physics; SPARC database; falsification criteria.
